% ---------------------------------------------------------------------------
% (a) The one Feynman diagram of q qbar -> b bbar at leading order, with the
% factors that build the matrix element; (b) the kinematics in the partonic
% centre-of-mass frame and the pTHat cut as a tube around the beam. Numbers
% from section [2] of the -v log of standalone_qqbar_bbbar_xsec.py.
% ---------------------------------------------------------------------------
\begin{tikzpicture}[x=1cm,y=1cm,
  every node/.style={font=\footnotesize},
  fermion/.style={thick,postaction={decorate},decoration={markings,mark=at position 0.55 with {\arrow[scale=1.1]{Stealth}}}},
  antifermion/.style={thick,postaction={decorate},decoration={markings,mark=at position 0.45 with {\arrowreversed[scale=1.1]{Stealth}}}},
  gluon/.style={thick,decorate,decoration={coil,aspect=0,segment length=4.5pt,amplitude=2pt}},
  vertex/.style={fill=black,circle,inner sep=1.6pt},
  lab/.style={font=\scriptsize,align=center},
  fac/.style={font=\scriptsize,align=center,text=orange!80!black},
  head/.style={font=\small\bfseries},
  beam/.style={thick,gray!70},
  tube/.style={dashed,red!70!black,thick},
  mom/.style={-{Stealth[length=2.5mm,width=2mm]},very thick}]
% ---------------- (a) the diagram ------------------------------------------------------------------
\node[head,anchor=west] at (-0.3,3.2) {(a) the only diagram: an $s$-channel gluon};
\coordinate (V1) at (1.6,0);
\coordinate (V2) at (4.2,0);
\draw[fermion] (-0.2,1.6) -- (V1) node[pos=0.0,above left,lab] {$q(p_1)$, colour $i$};
\draw[antifermion] (-0.2,-1.6) -- (V1) node[pos=0.0,below left,lab] {$\bar q(p_2)$, colour $j$};
\draw[gluon] (V1) -- (V2);
\draw[fermion] (V2) -- (6.0,1.6) node[above right,lab] {$b(p_3)$, colour $k$, mass $\mb$};
\draw[antifermion] (V2) -- (6.0,-1.6) node[below right,lab] {$\bar b(p_4)$, colour $l$, mass $\mb$};
\node[vertex] at (V1) {}; \node[vertex] at (V2) {};
\node[lab,below] at (2.9,-0.25) {$g^*$, $q^2=(p_1+p_2)^2=\shat$};
\node[fac,below] at (2.9,-0.7) {propagator $1/\shat$};
\node[fac,anchor=south] at (1.6,0.35) {$g_{\mathrm s}\,T^a_{ji}$};
\node[fac,anchor=south] at (4.2,0.35) {$g_{\mathrm s}\,T^a_{kl}$};
\node[lab,text=black!65,text width=6.6cm,anchor=north west] at (-0.3,-2.1) {$|\mathcal M|^2\propto g_{\mathrm s}^4=16\pi^2\alphas^2$; colour: $\tfrac19\sum|T^a_{ji}T^a_{kl}|^2=\tfrac19\cdot\mathrm{Tr}(T^aT^b)\mathrm{Tr}(T^aT^b)=\tfrac29$; spin average $\tfrac14$; angular structure $1+\cos^2\theta$ from a spin-1 exchange between two spin-$\tfrac12$ pairs. The result is \cref{eq:ancestral_home:dsigma}.};
% ---------------- (b) the kinematics in the partonic frame -----------------------------------------
\begin{scope}[xshift=9.4cm]
\node[head,anchor=west] at (-0.3,3.2) {(b) the partonic frame and the $\pthat\ge100\GeV$ tube};
\draw[beam] (-0.3,0) -- (7.3,0) node[right,lab] {beam ($z$)};
\draw[tube] (-0.3,1.3) -- (7.3,1.3) node[right,lab,text=red!70!black] {$\pthat^{\min}$};
\draw[tube] (-0.3,-1.3) -- (7.3,-1.3);
\coordinate (O) at (3.5,0);
% incoming
\draw[mom,green!50!black] (1.0,0) -- (3.1,0) node[midway,above,lab] {$u$, $\sqrt{\shat}/2$};
\draw[mom,violet] (6.0,0) -- (3.9,0) node[midway,above,lab] {$\bar u$};
% outgoing b (red) and bbar (blue), angle theta* to +z; cos theta*(bbar) = -0.82
\coordinate (B) at ($(O)+(-2.1,1.47)$);
\coordinate (Bb) at ($(O)+(2.1,-1.47)$);
\draw[mom,blue!80!black] (O) -- (B); \node[anchor=south,lab,text=blue!80!black] at (B) {$\bar b$: $p^*=198\GeV$, $\cos\theta^*=-0.820$};
\draw[mom,red!80!black] (O) -- (Bb) node[below right,lab,text=red!80!black] {$b$: $p^*=198\GeV$};
\draw[-{Stealth},thin] (O) ++(0.9,0) arc[start angle=0,end angle=145,radius=0.9];
\node[lab] at ($(O)+(0.25,1.15)$) {$\theta^*=2.53$};
% pT as the projection
\draw[dashed,thin] (B) -- ($(O)+(-2.1,0)$);
\node[lab,anchor=east] at ($(O)+(-2.2,0.7)$) {$\pt=p^*\sin\theta^*$\\$=113\GeV$};
% the allowed cone: |cos theta| <= c_max
\fill[red!15,opacity=0.6] (O) -- ($(O)+(2.6,1.3)$) -- ($(O)+(2.6,0)$) -- cycle;
\fill[red!15,opacity=0.6] (O) -- ($(O)+(2.6,-1.3)$) -- ($(O)+(2.6,0)$) -- cycle;
\fill[red!15,opacity=0.6] (O) -- ($(O)+(-2.6,1.3)$) -- ($(O)+(-2.6,0)$) -- cycle;
\fill[red!15,opacity=0.6] (O) -- ($(O)+(-2.6,-1.3)$) -- ($(O)+(-2.6,0)$) -- cycle;
\node[lab,text=red!70!black] at ($(O)+(2.1,0.45)$) {forbidden:\\$|\cos\theta|>c_{\max}$};
\node[lab,text=black!65,text width=7.4cm,anchor=north west] at (-0.3,-1.7) {Both quarks must leave through the wall of the tube, never through its ends: $|\cos\theta|\le c_{\max}=\sqrt{1-(\pthat^{\min}/p^*)^2}=0.863$ at our $\shat$, which removes $14\%$ of the solid angle and $19\%$ of $\hat\sigma$. The whole frame is then boosted along $z$ by $Y=-2.59$; \pt does not change.};
\end{scope}
\end{tikzpicture}
